\decorate{后记}

\setstyle{后记}

\large
\begin{spacing}{1.5}
    笔尖游移至此，已是书的尾声。每次敲击键盘，思绪都如潮水般涌来，有太多太多的话想要对你说，
    却总觉得词不达意。特别是在挑选插图的时候，每一张照片都像一扇任意门，把我带回当时的场景——你的笑容，
    你的眼神，你说话时微微上扬的嘴角——这些细节总让我格外思念你，恨不得能穿越时空，给正在打字的自己一个拥抱。

    说实话，写作的过程中我常常感到自己的文笔太过笨拙。这里的每一个字，都是我反复斟酌后留下的；
    每一张配图，都是我精心挑选、调整、排版的。没有参考任何人的建议，因为这份心意，我必须亲自完成。
    有时候为了一页的排版，我会反复调整到深夜，直到屏幕上的文字和图片呈现出最接近我想表达的样子。

    我知道这算不上什么珍贵的礼物，甚至可能有些幼稚。我们共同经历的点点滴滴，
    怎么可能被这区区一百多页的纸张完全承载？那些深夜的长谈，那些相视而笑的瞬间，那些只有我们懂的默契，
    都远远超出了文字能够表达的范围。我之所以执意要完成这本书，不过是想要用一种具体的方式，
    证明我们曾经这样真实地爱过、笑过、存在过。

    也许当你收到这本书时，会因为它朴素的装帧而失望；也许你会把它放在书架最不起眼的角落，任时光落满灰尘；
    又或许，你会珍而重之地将它收藏，偶尔翻开，重温我们共同走过的岁月。无论如何，我都理解并尊重你的选择。
    因为这本书既属于你，也属于我，它承载的是我全部的真挚，而你有权以任何方式安放这份心意。

    在整理这些文字时，我常常觉得，遇见你就像是遇见了世界上的另一个自己。我习惯性地自卑，总觉得自己不够好，
    不配拥有太多的善意；我们都过分敏感，常常因为一句无心的话语就陷入自我怀疑；我们都小心翼翼地活着，
    想要勇敢去爱却又害怕受伤。我们总是说要对自己好一点，可往往最先辜负的，就是自己。

    但正是通过爱你，通过这段旅程，我才真正开始学会接纳自己。在路上，我看见了内心那些忽明忽暗的理想，
    直面了贯穿成长历程的恐惧与怯懦，并在每一次的前行中将它们慢慢瓦解。所以，亲爱的，
    以后当你面临“要不要去”的抉择时，请勇敢地上路吧。不要犹豫，不要回头，这个世界有无数种可能，而唯有在路上，
    我们才能真正触达生命的各种维度。

    两周年快乐，我最爱的宝宝——涂琳。愿我们都能在爱与被爱中，成为更好的自己。

    \bottomtext{“我爱你，是我内心深处最深挚的告白”}
\end{spacing}

\clearpage


